\documentclass[11pt, a4paper]{article}
\usepackage[T1]{fontenc}
\usepackage[utf8]{inputenc}
\usepackage[a4paper, top=2.5cm, bottom=2.5cm, left=2.5cm, right=2.5cm]{geometry}
\usepackage{booktabs}
\usepackage{tabularx}
\usepackage{enumitem}
\usepackage[hidelinks]{hyperref}
\usepackage{parskip}
\usepackage{graphicx}
\usepackage{amsmath}
\usepackage{float}
\setlength{\parindent}{0pt}
\pagestyle{plain}

\begin{document}

% ------------------------------------------------------------------------------
% PORTADA
% ------------------------------------------------------------------------------
\begin{center}
{\Large \textbf{Pontificia Universidad Javeriana}}\\[0.3cm]
Facultad de Ciencias Economicas y Administrativas\\[0.3cm]
Estadistica -- 2026-II\\[1cm]
{\LARGE \textbf{La trampa del estrato: ?`el estrato socioeconomico determina el techo academico?}}\\[0.5cm]
{\large Informe Final del Proyecto}\\[1cm]
\textbf{Grupo 2}\\[0.3cm]
Carreno Cardenas Juan Jose\\
Gaitan Arias Santiago\\
Herrera Pulido Andres Felipe\\
Garavito Ortiz Isabela\\
Bonilla Martinez Mariana Alejandra\\[1cm]
27 de noviembre de 2026
\end{center}

\newpage
\tableofcontents
\newpage

% ------------------------------------------------------------------------------
\section{Introduccion}
% ------------------------------------------------------------------------------
% Contexto: estratificacion socioeconomica en Colombia y su relacion con
% resultados educativos.
% Pregunta: el estrato condiciona el desempeno en SBPro incluso controlando
% por habilidad previa (SB11)?
% Concepto de trampa de pobreza y reproduccion de desigualdad.

% ------------------------------------------------------------------------------
\section{Datos y Metodologia}
% ------------------------------------------------------------------------------
% Fuente: cruce Saber 11 -- Saber Pro (ICFES).
% Archivos: cruce\_saber11\_saberpro.parquet, saber11.parquet, saber\_pro.parquet.
% Variable dependiente: PUNT\_GLOBAL\_sbpro.
% Predictores: PUNT\_GLOBAL\_sb11, FAMI\_ESTRATOVIVIENDA\_sb11, INST\_ORIGEN\_sbpro,
%   FAMI\_EDUCACIONMADRE\_sb11, ESTU\_DEPTO\_RESIDE\_sbpro.
% Modelo 1 (aditivo), Modelo 2 (interaccion SB11 * estrato),
%   Modelo 3 (subsample carreras econ/admin).
% Justificacion de la interaccion: el efecto del estrato depende del nivel SB11?

% ------------------------------------------------------------------------------
\section{Analisis Exploratorio (Fase 1)}
% ------------------------------------------------------------------------------
% Distribucion de PUNT\_GLOBAL\_sbpro por estrato.
% Scatter SB11 vs SBPro coloreado por estrato.
% Estadisticas descriptivas por estrato y educacion de la madre.

% ------------------------------------------------------------------------------
\section{Intervalos de Confianza y Pruebas de Hipotesis (Fase 2)}
% ------------------------------------------------------------------------------
% IC para la media de PUNT\_GLOBAL\_sbpro por estrato.
% Prueba ANOVA o Kruskal-Wallis para diferencias entre estratos.
% Prueba chi-cuadrado para asociacion estrato e INST\_ORIGEN\_sbpro.

% ------------------------------------------------------------------------------
\section{Modelo de Regresion (Fase 3)}
% ------------------------------------------------------------------------------
% Modelo 1: aditivo (SB11 + estrato + inst\_origen + educ\_madre).
% Modelo 2: interaccion SB11 * estrato + inst\_origen + depto.
% Modelo 3: misma especificacion de M2 pero filtrado para carreras econ/admin.
% Tabla de coeficientes con IC al 95\%.
% Comparacion de modelos: R2 ajustado, AIC, BIC.
% Predicciones: estrato 1 vs estrato 5 al mismo percentil 50 de SB11.

% ------------------------------------------------------------------------------
\section{Diagnosticos y Robustez}
% ------------------------------------------------------------------------------
% Diagnosticos para cada modelo: residuales vs ajustados, QQ-plot.
% Test de Breusch-Pagan (homocedasticidad).
% VIF (multicolinealidad, incluyendo manejo de interacciones).
% Distancia de Cook (observaciones influyentes).
% Errores robustos HC1: comparacion con errores estandar clasicos.

% ------------------------------------------------------------------------------
\section{Resultados y Discusion}
% ------------------------------------------------------------------------------
% La interaccion es significativa? Que implica para la nocion de trampa?
% Comparacion de predicciones estrato 1 vs 5.
% Graficos de interaccion y coeficientes.
% Reflexion sobre reproduccion de desigualdad a traves del sistema educativo.

% ------------------------------------------------------------------------------
\section{Conclusiones}
% ------------------------------------------------------------------------------
% Respuesta a la pregunta de investigacion.
% Limitaciones del estudio.
% Implicaciones para politica de equidad y acceso a educacion superior.

% ------------------------------------------------------------------------------
\section*{Referencias}
% ------------------------------------------------------------------------------
% ICFES, literatura sobre estratificacion y desigualdad educativa en Colombia.

% ------------------------------------------------------------------------------
\section*{Anexo: Codigo R}
% ------------------------------------------------------------------------------
% Incluir codigo R utilizado (fase3\_grupo2.R) o referenciarlo.

\end{document}
